\documentclass{article}
\usepackage{hyperref}
\usepackage{Style}

\nocite{*} % Comentar si quiero citar
%\addbibresource{bibliografia.bib} % Quitar el comentado si quiero usar bibliografia

\begin{document}

\begin{minipage}{2.5cm}
    \includegraphics[width=2cm]{imagen_puc.jpg}
\end{minipage}
\begin{minipage}{12cm}
    {\sc Pontificia Universidad Católica de Chile\\
    Facultad de Matemáticas\\
    Departamento de Matemática\\
    Professor: Jesse Huang -- Student: Benjamín Mateluna}
\end{minipage}
\vspace{2ex}

{\centerline{\bf Selected Topics on Complex Algebraic Varieties - MPG3320 III}
\centerline{\bf Assignment 1}}
\centerline{\bf August 31, 2026}

\vspace{1cm}
\noindent\textbf{Problem 1.} Como el ideal $(x,y)$ es maximal, basta probar que
\begin{equation*}
    I=(x+xy,y+xy,x^{2},y^{2})\supseteq(x,y)
\end{equation*}
En efecto, notemos lo siguiente
\begin{equation*}
    (x+xy)(y+xy)=xy+x^{2}y+xy^{2}+x^{2}y^{2}
\end{equation*}
como $x^{2},y^{2}\in I$ se tiene que $xy\in I$, lo que implica que $x,y\in I$ y por lo tanto 
$(x,y)\subseteq I$.

\vspace{5mm}
\noindent\textbf{Problem 2.}
\begin{enumerate}
    \item[(a)] El conjunto $S$ no es una variedad afín, de lo contrario, sería un cerrado en
    la topología Zariski y, por consiguiente, un cerrado bajo la topología euclideana, lo que es 
    una contradicción.

    \item[(b)] Consideremos el ideal
    \begin{equation*}
        I=(x^{p-1}-1,x-y)\subseteq\mathbb{F}_{q}[x,y]
    \end{equation*}
    Dado $(a,a)\in S$, por pequeño teorema de fermat resulta que
    \begin{equation*}
        a^{p-1}-1=0 \hhtext{y} a-a=0
    \end{equation*}
    se sigue que $(a,a)\in V(I)$. Por otro lado, si $(x,y)\in V(I)$ vemos que $x=y$ y $x^{p-1}=1$,
    es decir, $x\neq0$ y así $(x,y)\in S$.

    \item[(c)] Definimos
    \begin{equation*}
        \U:=U_{0}\cap U_{1}\subseteq\mathbb{P}_{\C}^{1} \hhtext{y} 
        \V=\overline{V(xy-1)}\cap U_{2}\subseteq\mathbb{P}_{\C}^{2}
    \end{equation*}
    donde $U_{i}=\{[x_{0}:\cdots:x_{n}]\in\mathbb{P}_{\C}^{n}:x_{i}\neq0\}$. Puesto que el ideal
    $(xy-1)$ esta generado por un elemento su homogenización es $(xy-z^{2})$, luego
    $\V=V^{p}(xy-z^{2})\subseteq U_{i}$ para $i=0,1,2$.

    \vspace{1mm}
    Notar que tanto $\U$ como $\V$ son cuasi variedades proyectivas, más aún, mediante la 
    identificación natural, vemos que $\C^{*}$ es homeomorfo a $\U$ y $V(xy-1)$ con $\V$. Por
    otro lado, la aplicación
    \begin{equation*}
        z\mapsto(z,z^{-1})
    \end{equation*}
    es una correspondencia entre los puntos de $\C^{*}$ y $\V(xy-1)$, pero no es aplicación 
    polinomial. Sin embargo, dado $[x:y]\in\U$ y bajo las identificaciones, tenemos lo siguiente
    \begin{equation*}
        [x:y]=[xy^{-1}:1]\mapsto xy^{-1}\mapsto(xy^{-1},x^{-1}y)\mapsto [xy^{-1}:x^{-1}y:1]
        =[x^{2}:y^{2}:xy]
    \end{equation*}
    El mapa $\varphi:\U\to\V$ dado por $[x:y]\mapsto[x^{2}:y^{2}:xy]$ es morfismo de cuasi 
    variedades proyectivas, ya que es polinomial en cada coordenada, más aún, es isomorfismo con 
    inversa $\psi([x:y:z])=[x:z]$. Observemos lo siguiente
    \begin{equation*}
        \psi\varphi([x:y])=\psi([x^{2}:y^{2}:xy])=[x^{2}:xy]=[x:y]
    \end{equation*}
    y
    \begin{equation*}
        \varphi\psi([x:y:z])=\varphi([x:z])=[x^{2}:z^{2}:xz]=[x^{2}:xy:xz]=[x:y:z]
    \end{equation*}
    Lo que prueba el isomorfismo de cuasi variedades proyectivas.
\end{enumerate}

\newpage
\noindent\textbf{Problem 3.}
\begin{itemize}
    \item[(a)] En primer lugar, probaremos el siguiente lema
    
    \vspace{1mm}
    \noindent\textbf{Lemma:} \emph{Sea $f\in\C[x_{1},\cdots,x_{n}]$ no nulo, entonces 
    $f^{-1}(0)$ tiene interior vacío.}

    \begin{proof}
        Procederemos por inducción en $n\in\N$. Si $n=1$, sabemos que $f$ tiene finitas raíces, en
        caso de tener interior no vacío tendría infinitas raíces, lo que prueba el caso base.
        Supongamos, por contradicción, que existe $U\subseteq\C^{n}$ un abierto tal que 
        $U\subseteq f^{-1}(0)$. Escribimos
        \begin{equation*}
            f=\sum_{k}f_{k}(x_{1},\cdots,x_{n-1})x_{n}^{k}
            \htext{con} f_{k}\in\C[x_{1},\cdots,x_{n-1}]
        \end{equation*}
        Sea $a=(a_{1},\cdots,a_{n})\in U$ y consideremos la función continua
        \begin{align*}
            i:\C &\to \C^{n} \\
            z & \mapsto (a_{1},\cdots,a_{n-1},z)
        \end{align*}
        Notemos que $i^{-1}(U)$ es un abierto no vacío de $\C$ y $g(z)=f(a_{1},\cdots,a_{n-1},z)$ 
        un polinomio que se anula en un conjunto abierto de $\C$, por el caso base, se tiene que
        $g\equiv0$. Luego, $f_{k}(a_{1},\cdots,a_{n-1})=0$ para todo $k$.

        \vspace{1mm}
        El argumento anterior vale para todo $a\in U$ y como el mapa 
        $a\mapsto(a_{1},\cdots,a_{n-1})$ es abierto, el polinomio $f_{k}$ se anula en un conjunto
        abierto y por hipotesis inductiva $f_{k}\equiv0$ para todo $k$. Concluimos que $f^{-1}(0)$ 
        tiene interior vacío.
    \end{proof}

    El lema implica inmediatamente que dado $f\in\C[x_{1},\cdots,x_{n}]$ no constante, entonces
    $f^{-1}(\C\setminus\{0\})$ es denso en $\C^{n}$.

    \vspace{1mm}
    Sea $V$ una variedad algebraica afín propia, entonces $V=V(I)$ con 
    $I\subseteq\C[x_{1},\cdots,x_{n}]$ un ideal no nulo. Por el teorema de la base de Hilbert 
    existen $f_{1},\cdots,f_{m}$ tales que $I=(f_{1},\cdots,f_{m})$, así
    \begin{equation*}
        V=V(I)=V(f_{1},\cdots,f_{m})=\bigcap_{i=1}^{m}V(f_{i})\subseteq V(f_{1})
    \end{equation*}
    Tomando complemento y luego clausura, se obtiene que $V^{c}$ es denso en $\C^{n}$, es decir, 
    todo abierto en la topología Zariski es denso según la topología euclideana.

    \item[(b)] Supongamos, por contradicción, que ambas topologias son iguales. Consideremos
    $V=V(x-y)\subseteq\A^{2}_{\C}$ un cerrado según la topología Zariski y por lo tanto un cerrado 
    según la topología producto.

    \vspace{1mm}
    Lo anterior equivale a que la topología Zariski en $\A^{1}_{\C}$ es Hausdorff, pero es sabido 
    que dicha topología es idéntica a la topología cofinita en $\C$, la que no es Hausdorff en 
    conjuntos infinitos. Esta contradicción prueba la afirmación.
\end{itemize}

\vspace{5mm}
\noindent\textbf{Problem 4.}
\begin{enumerate}
    \item[(a)] Vemos que el ideal
    \begin{equation*}
        I:=(x_{1}x_{2})\subseteq k[x_{1},\cdots,x_{n}]
    \end{equation*}
    no es primo, en efecto, notar que $x_{1}x_{2}\in I$, pero $x_{i}\not\in I$. No obstante,
    resulta ser un ideal radical.

    \vspace{1mm}
    Si $f^{n}\in I$, por definición tenemos que $x_{1}x_{2}|f^{n}$, en particular, $x_{i}|f^{n}$ y 
    como este último es primo, se sigue que $x_{i}|f$. Dado que $x_{1}$ y $x_{2}$ son coprimos 
    se tiene que $x_{1}x_{2}|f$, lo que concluye la demostración.

    \item[(b)]
    \begin{itemize}
        \item $\Rightarrow|$ Supongamos que $I(V)$ no es primo, entonces existen $f_{1},f_{2}$ 
        polinomios tales que $f_{1}f_{2}\in I(V)$, pero $f_{i}\not\in I(V)$. Notar que
        \begin{equation*}
            V=(V\cap V(f_{1}))\cup(V\cap V(f_{2}))
        \end{equation*}
        pero $(V\cap V(f_{i}))\neq V$, ya que existe $p_{i}\in V$ tal que $f_{i}(p_{i})\neq0$. 
        Luego, $V$ es reducible.

        \item $\Leftarrow|$ Supongamos que $V$ es reducible, entonces
        \begin{equation*}
            V=V_{1}\cup V_{2}
        \end{equation*}
        con $V_{i}\neq V$ variedades afines. Entonces existen $p_{i}\in V\setminus V_{i}$ y 
        $f_{i}$ un polinomio tales que $f_{i}|_{V_{i}}\equiv0$ y $f_{i}(p_{i})\neq0$. Notemos que 
        $f_{1}f_{2}\in I(V)$, pero $f_{i}\not\in I(V)$. Concluimos que $I(V)$ no es primo.
    \end{itemize}
\end{enumerate}

\newpage
\noindent\textbf{Problem 5.} Veamos el siguiente lema

\vspace{1mm}
\noindent\textbf{Lemma:} \emph{Sea $k$ no algebraicamente cerrado, entonces existe un polinomio 
$q\in k[x,y]$ tal que $q(x,y)=0$ si y solo si $x=y=0$.}

\begin{proof}
    Como $k$ no es algebraicamente cerrado, existe $f\in k[x]$ un polinomio no constante sin 
    raíces en $k$. Escribamos
    \begin{equation*}
        f=a_{0}+a_{1}x+\cdots+a_{n}x^{n} \htext{con}a_{0},a_{n}\neq0
    \end{equation*}
    Definimos $q(x,y)\in k[x,y]$ como
    \begin{equation*}
        q(x,y):=\widetilde{f}=a_{0}y^{n}+a_{1}xy^{n-1}+\cdots+a_{n-1}x^{n-1}y+a_{n}x^{n}
    \end{equation*}
    Sea $(x,y)\in k^{2}$ tal que $q(x,y)=0$. Supongamos que $y\neq0$, de lo contrario, 
    $a_{n}x^{n}=0$ y entonces $x=0$. Luego,
    \begin{equation*}
        0=a_{0}+a_{1}\left(\frac{x}{y}\right)+\cdots+a_{n}\left(\frac{x}{y}\right)^{n}
        =f\left(\frac{x}{y}\right)
    \end{equation*}
    Pero la ecuación anterior no tiene solución en $k$ por la elección de $f$. Concluimos que
    $q(x,y)=0$ si y solo si $x=y=0$.
\end{proof}

\vspace{1mm}
\noindent Sea $V=V(f_{1},\cdots,f_{m})\subseteq\A_{k}^{n}$ una variedad algebraica, tomar 
$q\in k[x,y]$ como en el lema, definimos recursivamente los siguientes polinomios
\begin{equation*}
    q_{i}:=q(f_{i},q_{i+1}) \htext{para}1\leq i<m
\end{equation*}
y $q_{m}=f_{m}$. Afirmamos que $V(q_{1})=V$, en efecto, se tiene que $q_{1}=q(f_{1},q_{2})=0$ si y 
solo si $f_{1}=q_{2}=0$, del mismo modo, $q_{2}=0$ equivale a que $f_{2}=q_{3}=0$, inductivamente, 
tenemos que $q_{1}=0$ si y solo si $f_{i}=0$ para todo $i$, lo que prueba lo pedido.

\vspace{5mm}
\noindent\textbf{Problem 6.}
\begin{enumerate}
    \item[(a)] Sea $p\in C(Y)$ no nulo, por definición $[p]\in Y$, entonces $f(p)=0$ para todo 
    $f\in I$ homogéneo. Sea $f\in I$, escribimos
    \begin{equation*}
        f=\sum_{r}f_{r} \htext{con $f_{r}$ homogéneo}
    \end{equation*}
    Como $I$ es ideal homogéneo $f_{r}\in I$, entonces $f_{r}(p)=0$ y por lo tanto $f(p)=0$, luego
    $p\in V(I)$.

    \vspace{1mm}
    Sea $p\in V(I)$ no nulo, por lo tanto $f(p)=0$ para todo $f\in I$, en particular, para todo 
    $f\in I$ homogéneo y por ende $[p]\in Y$, es decir, $p\in C(Y)$.

    \item[(b)] Similar al caso de variedades afines, se verifica que una variedad proyectiva $V$ 
    es irreducible si y solo si $I^{h}(V)$ es un ideal primo. De este modo, basta probar que
    $I(C(Y))=I^{h}(Y)$.

    Sea $f\in I(C(Y))$, entonces $f|_{C(Y)}\equiv0$. Digamos que
    \begin{equation*}
        f=\sum_{r}f_{r} \htext{con $f_{r}$ homogéneo}
    \end{equation*}
    Queremos probar que $f_{r}|_{Y}\equiv0$ para todo $r$. Dado $[p]\in Y$, notamos que 
    $p\in C(Y)$, lo que implica que $f(p)=0$, más aún, $f(\lambda p)=0$ para todo $\lambda\neq0$.
    Entonces el polinomio
    \begin{equation*}
        g(\lambda)=f(\lambda p)=\sum_{r}f_{r}(p)\lambda^{r}
    \end{equation*}
    tiene infinitas raíces, lo que implica que $g$ es el polinomio nulo. Se sigue que $f_{r}(p)=0$
    para todo $r$. Concluimos que $f\in I^{h}(Y)$.

    Dado $f\in I^{h}(Y)$, digamos que
    \begin{equation*}
        f=\sum_{r}f_{r} \htext{con $f_{r}$ homogéneo y}f_{r}|_{Y}\equiv0
    \end{equation*}
    Fijemos $p\in C(Y)$ no nulo, nuevamente por definición tenemos que $[p]\in Y$ y se tiene que 
    $f_{r}(p)=0$ para todo $r$, por lo tanto $f(p)=0$. En resumen, $f\in I(C(Y))$, lo que prueba
    lo pedido.
\end{enumerate}

%\printbibliography % Quitar el comentado si quiero usar bibliografia

\end{document}
